\begin{exercise}[Identity transition]\label{ex:vi20-01}
Explain why setting $U_{ii}=X_i$ and $\varphi_{ii}=\operatorname{id}$ makes the gluing conditions uniform.
\end{exercise}

\begin{exercise}[Symmetry of the relation]\label{ex:vi20-02}
Show $\varphi_{ji}=\varphi_{ij}^{-1}$ implies the point-identification relation is symmetric.
\end{exercise}

\begin{exercise}[Transitivity]\label{ex:vi20-03}
Use the cocycle law to prove transitivity of the gluing relation.
\end{exercise}

\begin{exercise}[Open chart image]\label{ex:vi20-04}
Prove $\psi_i(X_i)$ is open in the quotient topology.
\end{exercise}

\begin{exercise}[Chart injectivity]\label{ex:vi20-05}
Prove $\psi_i:X_i\to X$ is injective.
\end{exercise}

\begin{exercise}[Overlap identity]\label{ex:vi20-06}
Show $\psi_i(X_i)\cap\psi_j(X_j)=\psi_i(U_{ij})$.
\end{exercise}

\begin{exercise}[Componentwise restrictions]\label{ex:vi20-07}
Show the compatible-family definition of $\OO_X$ is stable under restriction.
\end{exercise}

\begin{exercise}[Sheaf locality]\label{ex:vi20-08}
Prove uniqueness in the sheaf axiom for the glued presheaf.
\end{exercise}

\begin{exercise}[Sheaf gluing]\label{ex:vi20-09}
Prove existence in the sheaf axiom for $\OO_X$.
\end{exercise}

\begin{exercise}[Stalk identification]\label{ex:vi20-10}
For $x=\psi_i(x_i)$, prove $\OO_{X,x}\cong\OO_{X_i,x_i}$.
\end{exercise}

\begin{exercise}[Locality]\label{ex:vi20-11}
Why is the glued ringed space locally ringed?
\end{exercise}

\begin{exercise}[Scheme condition]\label{ex:vi20-12}
Why is the glued locally ringed space a scheme?
\end{exercise}

\begin{exercise}[Two-chart global sections]\label{ex:vi20-13}
For $X=U\cup V$, derive the equalizer formula for $\Gamma(X,\OO_X)$.
\end{exercise}

\begin{exercise}[Doubled origin sections]\label{ex:vi20-14}
Compute global sections of the doubled-origin line.
\end{exercise}

\begin{exercise}[Doubled origin non-affine]\label{ex:vi20-15}
Use the previous exercise to show the doubled-origin line is not affine.
\end{exercise}

\begin{exercise}[Inversion overlap]\label{ex:vi20-16}
Write the ring isomorphism used to glue two affine lines by inversion.
\end{exercise}

\begin{exercise}[Common distinguished refinement]\label{ex:vi20-17}
Let $x\in U\cap V$ for affine opens $U,V$. Outline the shrinking argument producing $W=D(f)\subseteq U=D(g)\subseteq V$.
\end{exercise}

\begin{exercise}[Common overlap ring]\label{ex:vi20-18}
If $W=D(f)\subseteq\Spec A=D(g)\subseteq\Spec B$, identify its coordinate ring in two ways.
\end{exercise}

\begin{exercise}[Compatible maps]\label{ex:vi20-19}
Given $f_i:X_i\to Y$ agreeing on overlaps, explain how the point maps glue.
\end{exercise}

\begin{exercise}[Uniqueness of a glued map]\label{ex:vi20-20}
Why is a morphism $X\to Y$ determined by its restrictions to the chart cover?
\end{exercise}

\begin{exercise}[Reconstruct from affine cover]\label{ex:vi20-21}
Given $X=\bigcup U_i$, specify the transition maps that reconstruct $X$ by gluing.
\end{exercise}

\begin{exercise}[Disjoint union as gluing]\label{ex:vi20-22}
Recover the disjoint union construction from the gluing theorem.
\end{exercise}

\begin{exercise}[Full-overlap gluing]\label{ex:vi20-23}
If $X_1\cong X_2$ and the full schemes are glued by that isomorphism, identify the result.
\end{exercise}

\begin{exercise}[Triple cocycle]\label{ex:vi20-24}
State precisely what can fail if $\varphi_{13}\ne\varphi_{23}\circ\varphi_{12}$ on a triple overlap.
\end{exercise}

\begin{exercise}[Inversion gluing global sections]\label{ex:vi20-25}
Show that the two affine lines glued by $u=t^{-1}$ have only constant global functions.
\end{exercise}

\begin{exercise}[Boundary audit]\label{ex:vi20-26}
Which numbered AG3 problem remains reserved immediately after VI/20, and where is it assigned?
\end{exercise}

